\documentclass[11pt,a4paper]{article}

\usepackage[utf8]{inputenc}
\usepackage[spanish, es-tabla]{babel}
\usepackage{geometry}
\usepackage{amsmath, amssymb}
\usepackage{tabularx}
\usepackage{booktabs}
\usepackage{xcolor} % Agregado para colores profesionales
\usepackage{tikz}
% Librerías modernas: positioning para ubicación relativa, babel para conflictos con español, arrows.meta para flechas Stealth
\usetikzlibrary{positioning, arrows.meta, babel} 
\usepackage{titlesec}
\usepackage{enumitem}
\usepackage{hyperref} % hyperref siempre debe ser el último paquete de formato

\geometry{top=2.5cm, bottom=2.5cm, left=2.5cm, right=2.5cm}
\hypersetup{colorlinks=true, linkcolor=blue!70!black, urlcolor=blue!70!black}

% Formato de secciones
\titleformat{\section}{\large\bfseries\uppercase}{\thesection.}{1em}{}
\titleformat{\subsection}{\normalsize\bfseries}{\thesubsection}{1em}{}

\begin{document}

\begin{center}
    \textbf{\Large GUÍA OFICIAL DEL ESTUDIANTE: PROYECTO FINAL INTEGRADOR (PFI)} \\
    \vspace{0.3cm}
    \textbf{\large UNIVERSIDAD CATÓLICA BOLIVIANA ``SAN PABLO'' -- SEDE TARIJA} \\
    \vspace{0.2cm}
    Departamento: Ciencias Exactas e Ingenierías | Carrera: Ingeniería Mecatrónica \\
    Asignatura: Robótica (IMT-342) | Gestión: Semestre II-2026 \\
    Docente: M.Sc. Ing. Bernardo Quiroga Turdera \\
    \vspace{0.4cm}
    \textit{\textbf{Documento: Pliego de Especificaciones Técnicas y Contrato de Ingeniería (PFI)}}
\end{center}

\vspace{0.5cm}
\section{Ficha Técnica y Objetivos SMART}

\subsection*{Título Oficial del Proyecto}
\textbf{Celda de Manufactura Flexible Inteligente: Clasificación y Paletizado Autónomo por Visión Artificial sobre el ABB IRB 120.}

\subsection*{Definición del Problema de Ingeniería}
En la industria moderna de manufactura y distribución logística, la variabilidad de los productos exige celdas de automatización flexible capaces de adaptarse a cambios sin necesidad de reprogramación rígida ni mecanismos de guiado mecánico. 

El equipo de ingeniería debe concebir, diseñar, implementar y operar (marco CDIO) una estación mecatrónica autónoma. Un operario colocará manualmente una pieza industrial (prismática o cilíndrica) en una posición $(x, y)$ y orientación $(\theta)$ arbitraria e impredecible dentro de una mesa de trabajo. El sistema deberá:
\begin{itemize}[label=\checkmark]
    \item Capturar y rectificar la imagen mediante una cámara USB cenital.
    \item Detectar la pieza y calcular su matriz de transformación espacial respecto a la base del robot empleando homografía plana resuelta por Descomposición en Valores Singulares (SVD).
    \item Generar una trayectoria continua con perfiles de velocidad suaves libres de singularidades.
    \item Transmitir las órdenes por red industrial al controlador ABB IRC5 para aprehender la pieza con una pinza neumática y paletizarla ordenadamente.
\end{itemize}

\vspace{0.4cm}
\begin{center}
\begin{tikzpicture}[
    node distance=0.6cm, 
    auto, 
    thick,
    box/.style={draw=blue!80!black, fill=blue!5, rectangle, rounded corners, minimum width=7cm, minimum height=0.8cm, align=center, font=\small\bfseries},
    arrow/.style={->, >=Stealth, thick, draw=gray!80!black}
]
    
    \node[box] (cam) {[Cámara USB Cenital (2D)]};
    \node[box, below=of cam] (cv) {[Nodo OpenCV / ROS (EC3)]};
    \node[box, below=of cv] (ik) {[Resolvedor Pieper (EC2)]};
    \node[box, below=of ik] (ros) {[Driver ROS-Industrial (EC4)]};
    \node[box, below=of ros] (irc5) {[Controlador ABB IRC5]};
    \node[box, below=of irc5] (robot) {[Robot ABB IRB 120 Real]};
    
    \draw[arrow] (cam) -- node[right, font=\footnotesize] {Flujo de Video} (cv);
    \draw[arrow] (cv) -- node[right, font=\footnotesize] {Pose en mm} (ik);
    \draw[arrow] (ik) -- node[right, font=\footnotesize] {Trayectoria LSPB} (ros);
    \draw[arrow] (ros) -- node[right, font=\footnotesize] {Socket TCP/IP} (irc5);
    \draw[arrow] (irc5) -- node[right, font=\footnotesize] {Comandos RAPID} (robot);
\end{tikzpicture}
\end{center}

\subsection*{Objetivos Cuantitativos del Proyecto (KPIs de Desempeño)}
\begin{itemize}
    \item \textbf{KPI-1 (Precisión Posicional):} Error absoluto lineal de posicionamiento final entre el centro de la pinza y el centroide real de la pieza $\Delta E < 1.5\text{ mm}$.
    \item \textbf{KPI-2 (Eficiencia Temporal):} Tiempo de ciclo total $t_{\text{ciclo}} < 8\text{ s}$ por pieza (desde detección hasta descarga).
    \item \textbf{KPI-3 (Tasa de Éxito):} $> 90\%$ de clasificaciones autónomas correctas en 10 ejecuciones continuas sin colisiones.
    \item \textbf{KPI-4 (Seguridad Industrial):} Detención automática de emergencia ante desconexión de red con latencia de respuesta $< 15\text{ ms}$ cumpliendo la norma ISO 10218-1/2.
\end{itemize}

\section{Arquitectura de Hardware y Software (BOM)}
Para garantizar la homogeneidad en el desarrollo, cada equipo utilizará exclusivamente la siguiente infraestructura autorizada:

\vspace{0.3cm}
\noindent
\begin{tabularx}{\textwidth}{@{} l l X @{}}
\toprule
\textbf{Categoría} & \textbf{Componente} & \textbf{Especificación Autorizada} \\
\midrule
Hardware & Manipulador Industrial & ABB IRB 120 de 6 GDL (Alcance: $580\text{ mm}$, Carga: $3\text{ kg}$). \\
Controlador & Unidad de Potencia & ABB IRC5 con módulo Ethernet TCP/IP. \\
Efector Final & Actuador Neumático & Pinza de doble efecto operada por electroválvulas (24 V). \\
Sensor Óptico & Cámara USB Cenital & Sensor CMOS HD/FullHD, enfoque manual fijado, montaje a $80\text{ cm}$. \\
Sistema Operat. & Entorno de Desarrollo & Linux Ubuntu 24.04 LTS (nativa o dual-boot). \\
Middleware & Ecosistema Robótico & ROS (Robot Operating System) versión Noetic / Humble. \\
Librerías & Entorno de Cómputo & Python 3.x, NumPy, OpenCV. \\
Industrial & Simulación y Lenguaje & ABB RobotStudio y lenguaje nativo RAPID \texttt{.mod}. \\
Metrología & Instrumento Verificación & Calibrador Vernier digital de precisión ($\pm 0.02\text{ mm}$). \\
\bottomrule
\end{tabularx}

\section{Condiciones de Frontera y Reglas ``Anti-Caja Negra''}
Para certificar la apropiación de las competencias matemáticas, el proyecto está sujeto a estas restricciones:
\begin{enumerate}
    \item \textbf{Prohibición de MoveIt! o Resolvedores Automáticos:} El cálculo de la cinemática inversa debe programarse en Python, aplicando analíticamente el Método de Pieper.
    \item \textbf{Prohibición de \texttt{cv2.findHomography} Directo:} La estimación de la matriz de homografía $H$ debe programarse construyendo manualmente la matriz $A \in \mathbb{R}^{2N \times 9}$ del algoritmo DLT y extrayendo el espacio nulo mediante SVD.
    \item \textbf{Prohibición de Coordenadas ``Hardcodeadas'':} Las posiciones de picking no pueden estar pregrabadas. Si el robot se desplaza sin que la cámara publique la pose en ROS, la prueba será anulada.
    \item \textbf{Desactivación Estricta del Autofocus:} El enfoque de la cámara debe fijarse manualmente para evitar variar la distancia focal ($f_x, f_y$), lo que invalidaría la calibración intrínseca $K$.
\end{enumerate}

\section{Desglose Estructurado de Hitos (WBS) y Pasaportes}
El proyecto se dividirá en cuatro (4) hitos evaluativos vinculados a las Unidades de la materia.

\begin{description}[style=nextline, leftmargin=*]
    \item[\textbf{Hito 1: Gemelo Digital URDF y Cinemática Directa (Semana 5)}]
    \textbf{Entregable:} Paquete de ROS con el URDF modular del ABB IRB 120 acoplado con la pinza. \\
    \textbf{Pasaporte de Aprobación:} La pose cartesiana $[x, y, z]$ de RViz debe coincidir con la función analítica en Python, registrando un error $\Delta E < 10^{-6}\text{ m}$.

    \item[\textbf{Hito 2: Motor Cinemático Inverso (Pieper) y Perfiles LSPB (Semana 9)}]
    \textbf{Entregable:} Nodo de ROS en Python que reciba una pose, ejecute Pieper y genere una trayectoria LSPB. \\
    \textbf{Pasaporte de Aprobación:} Demostración en Jupyter Notebook de la divergencia analítica de velocidades y torques en singularidad ($q_5 = 0\text{ rad}$).

    \item[\textbf{Hito 3: Localización Hand-Eye SVD y Segmentación en OpenCV (Semana 13)}]
    \textbf{Entregable:} Archivo \texttt{camera\_info.yaml} y nodo de visión que publique la pose en milímetros. \\
    \textbf{Pasaporte de Aprobación:} Error de reproyección intrínseca $\text{RMSE} < 0.5\text{ px}$ y error físico real $\Delta E < 1.5\text{ mm}$ medido con Vernier.

    \item[\textbf{Hito 4: Integración Física, Red ROS-Industrial y Artículo IEEE (Semana 16)}]
    \textbf{Entregable:} Operación física autónoma, código RAPID \texttt{.mod}, arquitectura FSM y artículo IEEE. \\
    \textbf{Pasaporte de Aprobación:} Clasificación autónoma de 5 piezas con $t_{\text{ciclo}} < 8\text{ s}$ por pieza y desconexión segura ante fallas de red.
\end{description}

\section{Matriz de Pruebas de Rendimiento y Medición del Error}
Durante las jornadas de defensa de los Hitos 3 y 4, se aplicará la siguiente metrología:

\vspace{0.3cm}
\begin{center}
\begin{tikzpicture}[
    node distance=0.5cm, 
    auto, 
    thick,
    box/.style={draw=red!80!black, fill=red!5, rectangle, rounded corners, minimum width=6.5cm, minimum height=0.6cm, align=center, font=\small},
    arrow/.style={->, >=Stealth, thick, draw=gray!80!black}
]
    
    \node[box] (p1) {[Pieza en Mesa (Coordenada Aleatoria)]};
    \node[box, below=of p1] (p2) {[Detección Óptica y Cálculo SVD]};
    \node[box, below=of p2] (p3) {[Descenso del Robot ABB IRB 120]};
    \node[box, below=of p3] (p4) {[Parada de Emergencia en Posición]};
    \node[box, below=of p4] (p5) {[Medición con Calibrador Vernier del Desvío $\Delta E$]};
    
    \draw[arrow] (p1) -- (p2);
    \draw[arrow] (p2) -- (p3);
    \draw[arrow] (p3) -- (p4);
    \draw[arrow] (p4) -- (p5);
\end{tikzpicture}
\end{center}

\vspace{0.3cm}
\noindent \textbf{Cálculo del Error Absoluto Posicional:}
$$ \Delta E = \sqrt{(x_{\text{calculado}} - x_{\text{real}})^2 + (y_{\text{calculado}} - y_{\text{real}})^2} $$
Si el error $\Delta E \ge 1.5\text{ mm}$, el intento se contabiliza como fallido.

\section{Estándares de Ingeniería de Software y Documentación}
\subsection*{Convención para Repositorios Git}
Se exige el uso de \textit{Conventional Commits} (ej. \texttt{feat(kinematics): add pieper solver class}). Las ramas principales serán \texttt{main}, \texttt{dev} y \texttt{feature/*}. Commits masivos al final del semestre serán penalizados.

\subsection*{Formato del Artículo Científico Final}
El informe final del Hito 4 debe redactarse obligatoriamente en formato de dos columnas IEEE (Transactions on Robotics) con una extensión de 4 a 6 páginas. Debe incluir: Abstract, Introduction, Kinematic \& Dynamic Modeling, Computer Vision, System Architecture, Experimental Results y Conclusions.

\section{Protocolo de Seguridad Industrial (ISO 10218-1/2)}
\begin{itemize}
    \item \textbf{Operación en Modo Manual Reducido:} Durante calibraciones (Hitos 3 y 4), el selector del IRC5 debe estar en T1 (velocidad al 25\%).
    \item \textbf{Pulsador de Habilitación (Deadman Switch):} Mantenerlo en posición intermedia; presionarlo a fondo o soltarlo debe detener el robot.
    \item \textbf{Zona de Exclusión:} Nadie debe ingresar al radio de $580\text{ mm}$ en modo automático.
    \item \textbf{Watchdog de Red:} El código RAPID debe detener el robot si la PC deja de enviar mensajes por más de $500\text{ ms}$.
\end{itemize}

\end{document}